% ============================================================
%  12-devops.tex
%  Chapitre 5 : DevOps et mise en production
%
%  Captures à placer dans rapport/overleaf/images/devops/ :
%    actions.png             -> GitHub > Actions (liste workflows)
%    docker-azure.png        -> SSH : docker compose ps
%    deploy-success.png      -> Actions > Deploy (run vert)
%    actuator-health.png     -> curl ou navigateur /actuator/health UP
%    grafana-dashboard.png   -> Grafana :3000 (si observability lancée)
%    vercel-login.png        -> https://solar-ease-sxyy.vercel.app/login OK
%  Schéma : devops.png -> Chaîne DevOps et déploiement
% ============================================================

\newcommand{\devopsfig}[4][0.88]{%
  \begin{figure}[H]
    \centering
    \IfFileExists{#2}{%
      \includegraphics[width=#1\textwidth]{#2}%
    }{%
      \fbox{\parbox[c][4cm][c]{#1\textwidth}{\centering
        \small\textit{Capture à ajouter}\\[4pt]
        \small\texttt{#2}\\[4pt]
        \small #3
      }}%
    }
    \caption{#4}
  \end{figure}
}

\chapter{DevOps et mise en production}

\begin{chapterintro}
Ce chapitre présente l'industrialisation de SolarEase~: containerisation,
CI/CD GitHub Actions, observabilité, sécurité, puis la \textbf{mise en
production réelle} (frontend Vercel, backend sur VM Azure). L'objectif est
de garantir des livraisons reproductibles, traçables et vérifiables~: chaque
\textit{push} déclenche des contrôles automatisés, le déploiement backend
s'effectue par SSH sans intervention manuelle sur la VM, et les figures
illustrent la chaîne DevOps ainsi que les preuves de bon fonctionnement en
production (captures d'écran).
\end{chapterintro}

% ══════════════════════════════════════════════════════════
\section{Architecture et principes}
% ══════════════════════════════════════════════════════════

La démarche s'appuie sur le cadre \textbf{CALMS}~:
\textbf{C}ulture (intégration continue partagée par l'équipe)~;
\textbf{A}utomation (workflows GitHub Actions)~;
\textbf{L}ean (pipelines courts, retour rapide)~;
\textbf{M}easurement (Prometheus, Grafana, health checks)~;
\textbf{S}haring (Docker Compose, manifests K8s, \texttt{RUNBOOK.md}).

Concrètement, le cycle de vie suit trois axes~: (1)~développement local
via Vite + Compose~; (2)~intégration continue à chaque \textit{push} ou PR
(build Maven, tests, scans)~; (3)~déploiement de démonstration sur Azure
et Vercel, avec vérification post-déploiement (\texttt{/actuator/health}).

\devopsfig[0.95]{devops.png}{Ajouter \texttt{devops.png}.}{Chaîne DevOps et déploiement de démonstration}

% ══════════════════════════════════════════════════════════
\section{Containerisation}
% ══════════════════════════════════════════════════════════

Chaque microservice utilise un \texttt{Dockerfile} \textbf{multi-stage}~:
une étape \texttt{build} (Maven~3.9, Temurin~17) compile le JAR, une étape
\texttt{runtime} embarque uniquement le JAR sur une image JRE Alpine (ou
Debian pour \texttt{dimensioning-service}, avec Tesseract OCR pour l'extraction
de factures). Cette séparation réduit la taille des images et la surface
d'attaque.

Le fichier \texttt{backend/docker-compose.yml} orchestre l'ensemble sur le
réseau bridge \texttt{solarease-net}~:
\begin{itemize}[leftmargin=*, nosep]
  \item \textbf{Gateway} (8080)~: point d'entrée unique, routage JWT~;
  \item \textbf{Identity} (8081), \textbf{Project} (8082),
        \textbf{Dimensioning} (8083)~: microservices métier~;
  \item trois instances \textbf{PostgreSQL~16} (identity, project,
        dimensioning avec \texttt{pgvector})~;
  \item \textbf{Ollama} et \textbf{n8n} pour l'assistant IA et l'automatisation.
\end{itemize}

En local, le frontend Vite (\texttt{npm run dev}, port~5173) appelle
\texttt{/api} via le proxy Vite~; en production, c'est \texttt{vercel.json}
qui assure le même rôle (voir section~\ref{sec:prod-pfe}).

\devopsfig{docker-azure.png}{Terminal SSH sur la VM Azure~:
\texttt{cd \textasciitilde/SolarEase/backend \&\& docker compose ps}.}
{Stack Docker en production sur la VM Azure -- neuf conteneurs actifs}

% ══════════════════════════════════════════════════════════
\section{Pipeline CI/CD}
% ══════════════════════════════════════════════════════════

Quatre workflows GitHub Actions couvrent le cycle de livraison. Ils sont
déclenchés de façon \textbf{ciblée} (\texttt{paths}) afin d'éviter des runs
inutiles lors de modifications du rapport ou du frontend seul.

\begin{table}[H]
  \caption{Workflows GitHub Actions}
  \label{tab:workflows}
  \centering\small
  \begin{tabularx}{\textwidth}{|L{3.8cm}|L{4cm}|X|}
    \hline
    \rowcolor{black!12}
    \textbf{Fichier} & \textbf{Déclencheur} & \textbf{Rôle} \\
    \hline
    \texttt{backend-ci.yml} & push/PR \texttt{backend/**} &
    \texttt{mvn verify} en matrice (4 services), JaCoCo \\
    \hline
    \texttt{docker-publish.yml} & \texttt{main}, tags \texttt{v*} &
    Build et push images sur GHCR \\
    \hline
    \texttt{security-scan.yml} & push, PR, cron &
    Trivy (fs + images), Gitleaks \\
    \hline
    \texttt{deploy.yml} & push \texttt{main}, manuel &
    Déploiement VM Azure (SSH, \texttt{docker compose up}) \\
    \hline
  \end{tabularx}
\end{table}

Le workflow \texttt{deploy.yml} s'exécute via \texttt{appleboy/ssh-action}~:
connexion à la VM avec les secrets \texttt{AZURE} (IP), \texttt{AZURE\_VM}
(utilisateur) et \texttt{AZURE\_VM\_SSH\_KEY} (clé privée)~; puis
\texttt{git pull}, \texttt{docker compose up -d --build}, et boucle d'attente
sur \texttt{/actuator/health} (jusqu'à 3~minutes) avant de valider le déploiement.

\devopsfig{actions.png}{GitHub $\rightarrow$ SolarEase
$\rightarrow$ Actions. Montrer Backend CI, Deploy, Docker Publish,
Security Scan.}{Interface GitHub Actions -- workflows automatisés}

\devopsfig{devops2.png}{Actions $\rightarrow$ Deploy
$\rightarrow$ run réussi (coche verte).}{Workflow Deploy -- mise à jour
automatisée du backend sur Azure}

% ══════════════════════════════════════════════════════════
\section{Observabilité}
% ══════════════════════════════════════════════════════════

Chaque microservice Spring Boot expose \texttt{/actuator/health} (sonde de
vivacité) et \texttt{/actuator/prometheus} (métriques Micrometer). La stack
optionnelle est décrite dans
\texttt{backend/observability/docker-compose.observability.yml} et se lance
en complément du Compose principal~:

\begin{itemize}[leftmargin=*, nosep]
  \item \textbf{Prometheus} (9090)~: scrape toutes les 15~s les quatre
        services sur le réseau Docker~;
  \item \textbf{Grafana} (3000)~: dashboards provisionnés
        (\texttt{solarease.json}), sources Prometheus et Loki auto-configurées~;
  \item \textbf{Loki} (3100) + \textbf{Promtail}~: centralisation des logs
        conteneurs via le socket Docker.
\end{itemize}

Un correctif PFE a aligné les chemins \texttt{observability/} et la config
Loki (\texttt{ring.kvstore: inmemory}) pour un démarrage fiable sur VM à
ressources limitées. L'accès Grafana en production se fait via tunnel SSH
(port~3000 non exposé publiquement) ou en local après
\texttt{docker compose -f \ldots observability.yml up -d}.

\devopsfig{devops5.png}{Démarrer la stack
observability puis ouvrir \texttt{localhost:3000} ou l'IP VM :3000.}
{Dashboard Grafana -- métriques des microservices}

% ══════════════════════════════════════════════════════════
\section{Qualité et sécurité}
% ══════════════════════════════════════════════════════════

\textbf{JaCoCo} génère la couverture de code à l'étape \texttt{mvn verify}~;
\textbf{Trivy} analyse le système de fichiers et les images Docker~;
\textbf{Gitleaks} détecte les fuites de secrets dans l'historique Git.
Les résultats Trivy/Gitleaks sont publiés en \textbf{SARIF} pour consultation
dans l'onglet \textit{Security} de GitHub.

Les secrets applicatifs et d'infrastructure restent \textbf{hors Git}~:
fichier \texttt{backend/.env} (non versionné), secrets GitHub Actions pour
Azure et variables Vercel pour le frontend. Aucune clé SSH ni mot de passe
n'apparaît dans le dépôt~; la clé de déploiement est stockée chiffrée côté
GitHub et la clé publique correspondante dans
\texttt{\textasciitilde/.ssh/authorized\_keys} sur la VM.

% ══════════════════════════════════════════════════════════
\section{Mise en production (démo PFE)}
% ══════════════════════════════════════════════════════════
\label{sec:prod-pfe}

Le \textbf{backend} tourne sur une \textbf{VM Azure} (Ubuntu~24.04,
France Central, IP publique \texttt{4.233.29.236})~; le \textbf{frontend}
React/Vite est hébergé sur \textbf{Vercel}
(\url{https://solar-ease-sxyy.vercel.app}).

\paragraph{Déploiement backend.}
Clone du dépôt GitHub sur la VM (\texttt{\textasciitilde/SolarEase}),
\texttt{docker compose up -d --build}, NSG Azure avec la règle TCP~8080
ouverte. Le workflow \texttt{deploy.yml} automatise cette séquence à chaque
\textit{push} sur \texttt{main} modifiant \texttt{backend/**}.

\paragraph{Déploiement frontend.}
Vercel build le dossier frontend (\texttt{npm run build}, sortie \texttt{dist}).
Le fichier \texttt{vercel.json} réécrit \texttt{/api/:path*} et \texttt{/ws}
vers le gateway Azure, sans variable \texttt{VITE\_API\_URL} en production~:
le navigateur appelle le même domaine Vercel, qui proxifie vers la VM. Ainsi,
l'authentification JWT transite par \texttt{POST /api/auth/login} sur
\texttt{solar-ease-sxyy.vercel.app}, preuve que la chaîne frontend $\rightarrow$
rewrite $\rightarrow$ gateway $\rightarrow$ identity-service fonctionne de bout
en bout.

\devopsfig{devops8.png}{Navigateur ou terminal~:
\texttt{http://4.233.29.236:8080/actuator/health} avec statut UP.}
{Health check public du gateway Azure}

\devopsfig{devops10.png}{URL production Vercel, login
admin réussi (\texttt{admin@solarease.demo}).}{Application SolarEase en
production -- authentification via rewrite Vercel}

\begin{table}[H]
  \caption{Environnements SolarEase}
  \label{tab:environments}
  \centering\small
  \begin{tabularx}{\textwidth}{|L{2.2cm}|L{3.2cm}|L{3.5cm}|X|}
    \hline
    \rowcolor{black!12}
    \textbf{Env.} & \textbf{Backend} & \textbf{Frontend} & \textbf{Usage} \\
    \hline
    Local & Docker Compose & \texttt{npm run dev} & Développement \\
    \hline
    Production & VM Azure + Compose & Vercel & Démo PFE, URL publique \\
    \hline
  \end{tabularx}
\end{table}

% ══════════════════════════════════════════════════════════
\section{Infrastructure as Code (Kubernetes)}
% ══════════════════════════════════════════════════════════

Le dossier \texttt{k8s/} (Kustomize, overlays \texttt{dev}/\texttt{prod})
documente une montée en charge cloud-native~: Deployments par microservice,
Services ClusterIP, ConfigMaps et Secrets. Cette voie n'est pas activée pour
la démo PFE (coût et complexité), mais permet une évolution future via
\texttt{kubectl apply -k k8s/overlays/prod} sans réécrire les manifests.

% ══════════════════════════════════════════════════════════
\section{Tests de validation}
% ══════════════════════════════════════════════════════════

Les scénarios ci-dessous ont été exécutés manuellement ou via CI pour valider
la chaîne DevOps et la mise en production.

\begin{table}[H]
  \caption{Validation DevOps (extrait)}
  \label{tab:devops-tests}
  \centering\small
  \begin{tabularx}{\textwidth}{|C{1cm}|L{4.2cm}|X|C{1.2cm}|}
    \hline
    \rowcolor{black!12}
    \textbf{ID} & \textbf{Test} & \textbf{Résultat attendu} & \textbf{Statut} \\
    \hline
    D-01 & \texttt{docker compose up} & Stack UP $<$60s & \ok~OK \\
    \hline
    D-02 & PR sur backend & CI verte (4 services) & \ok~OK \\
    \hline
    D-03 & \texttt{deploy.yml} sur \texttt{main} & SSH + health UP & \ok~OK \\
    \hline
    D-07 & Logs $\rightarrow$ Loki & Filtre par conteneur & \ok~OK \\
    \hline
    D-13 & Health gateway Azure & \texttt{UP} sur :8080 & \ok~OK \\
    \hline
    D-14 & Login via Vercel & Auth OK via \texttt{/api} & \ok~OK \\
    \hline
  \end{tabularx}
\end{table}

\begin{conclusionbox}
SolarEase dispose d'une chaîne DevOps complète (build, tests, scans,
observabilité) et d'une mise en production opérationnelle Vercel + Azure.
Les captures d'écran du chapitre --- stack Docker sur VM, workflows Actions,
health check public, dashboard Grafana et login production --- constituent
les preuves visuelles de cette industrialisation et du critère D-14
(authentification bout en bout via rewrite).
\end{conclusionbox}
